\section{Results}
\label{sec:results}
\subsection{Replication of the industry benchmark}
Table~\ref{tab:replication} reports the complete benchmark comparison. The preliminary covariance correction raises the gross Sharpe ratio from 0.322 to 0.592. Further shrinkage raises it to 0.741 at $w=0.25$ and 0.816 at $w=0.75$. Annual selection based on preceding returns produces a Sharpe ratio of 0.770. The fixed-grid portfolios describe the response to shrinkage; the annual-selection portfolio provides the historically selected comparator. The replication uses annual trailing-window selection and monthly predicted-volatility scaling; Appendix~\ref{app:data} details their relationship to the published procedure. The central empirical pattern in \citet{pedersen2021} is recovered: regularization improves performance relative to the raw optimizer. The original and reproduced values are shown together, with the detailed discrepancy accounting deferred to Appendix~\ref{app:data}.
\input{generated/table_replication}

Figure~\ref{fig:pc} shows how regularization changes principal-component risk. In the smallest-eigenvalue group, the realized-to-predicted risk ratio falls from 5.22 under the raw covariance matrix to 1.96 under the preliminary correction and 0.94 under the selected EPO matrix. Under the common matrix $Q_t$ used for position sizing, that group's share of predicted portfolio variance falls from 82\% for raw mean--variance to 19\% for corrected mean--variance and 4\% for EPO. Realized component Sharpe ratios are higher in the largest-eigenvalue groups. Regularization therefore reduces concentration in directions with understated risk and shifts exposure toward directions where the signal performs better. The panels distinguish the three portfolio-specific risk matrices from the common position-sizing matrix.
\reportfigure{problem_pc.pdf}{Principal-component diagnostics for the replicated benchmark, 1942--2018. MVO denotes mean--variance optimization; OOS denotes selection using past data. Components are ordered by ascending eigenvalue and divided into five groups. The upper and lower-left panels use each portfolio's own risk matrix: $\widehat\Sigma_t$, $Q_t$, or $(1-w)Q_t+wV_t$. The lower-left panel takes component positions in the signal direction. The lower-right panel uses the common position-sizing matrix $Q_t$. All panels use the common 863-month sample. Appendix~\ref{app:math} defines the statistics.}{fig:pc}{1}

\subsection{Scalar compression and uncertainty-estimation error}
The scalar restriction costs approximately 12\% and 29\% of mean oracle utility in the profiles with actual coefficients of variation 54\% and 120\%. With a uniform profile, the heterogeneous allocation belongs to the scalar family and compression loss is essentially zero. The additional losses from choosing the mean scalar instead of the oracle scalar are approximately 0.5\% and 2.0\% of mean oracle utility. Thus most of the loss in the two nonuniform profiles remains even after the scalar parameter is optimized with knowledge of the uncertainty matrix.

Figure~\ref{fig:compression} separates these losses and relates compression to exposure alignment. Within the higher-dispersion profile, compression loss rises from 23\% of group mean oracle utility in the lowest alignment third to 36\% in the highest. The corresponding shares for the moderate profile are 10\% and 15\%. The profile remains fixed within each comparison; the sorting variable changes with the signal and its implied positions. Loss therefore depends on where portfolio exposure coincides with uncertainty differences. The loss shares summarize the relative cost of the scalar restriction in these known-input models. They describe the model's allocation trade-off and do not forecast a corresponding increase in historical returns. Appendix~\ref{app:scalar-sufficiency} characterizes when compression cost is small; Appendix~\ref{app:simulation} reports group means and Monte Carlo uncertainty in annual utility basis points.
\reportfigure{compression_loss.pdf}{Utility loss from scalar compression. The left panel separates compression loss after scalar optimization from the additional loss of the mean-scalar choice. The right panel sorts draws by exposure alignment within each profile. Each bar is mean loss divided by mean oracle utility within the displayed group. Error bars are 95\% Monte Carlo intervals for these ratios, accounting for the paired numerator and denominator. Each profile has 1,000 paired draws; intervals condition on that fixed profile. Target coefficients of variation 50\% and 100\% produce actual values 54\% and 120\%.}{fig:compression}{1}

The controlled estimation experiment compares these potential benefits with the loss caused by uncertainty-estimation error. At the specified mean scalar $\bar\kappa=2.5$, restriction costs are approximately 13\% and 31\% of mean oracle utility in the two nonuniform settings. These costs include the additional scalar-choice loss. With estimation-error log standard deviation $\delta=0.5$, mean $P$ is about 14\% of oracle utility in the moderate-heterogeneity setting and 10\% in the high-heterogeneity setting. The incremental utilities over the mean-scalar allocation are approximately $-1\%$ and $+21\%$ of mean oracle utility in the two settings.

In synthetic calibration, the bootstrap estimator has log-error dispersion $\delta_{\mathrm{cal}}=0.414$. At this measured dispersion, estimation losses are approximately 9.3\% and 6.7\% of oracle utility in the moderate- and high-heterogeneity settings, below restriction costs of 13\% and 31\%. The net gains are 3.7\% and 24.0\%, with 95\% Monte Carlo intervals of $[3.4,4.0]\%$ and $[23.4,24.6]\%$. Thus this error-dispersion match retains a positive gain in both nonuniform settings. At $\delta=1$, estimation loss exceeds restriction cost in both; mean net differences fall to approximately $-64\%$ and $-19\%$ of oracle utility. With uniform true uncertainty, the scalar is already optimal and estimation noise only reduces utility. Appendix~\ref{app:simulation} reports the full paired results.
\reportfigure{estimation_tradeoff.pdf}{The value of estimating asset-specific uncertainty. Each panel holds true uncertainty and signal draws fixed while varying the log standard deviation $\delta$ of multiplicative estimation error. Curves show mean restriction cost $L_{\bar\kappa}$ at $\bar\kappa=2.5$, estimation loss $P$, and incremental utility $L_{\bar\kappa}-P$, each divided by mean oracle utility. Error bars are 95\% paired Monte Carlo intervals for these ratios across 1,000 draws. The dashed line marks $\delta_{\mathrm{cal}}=0.414$ from zero-feedback estimator calibration; all curves are evaluated at this value. Labels at the dashed line give net gains as percentages of mean oracle utility. Lines connect evaluated error levels. At zero error the estimated heterogeneous allocation equals the oracle.}{fig:estimation-tradeoff}{1}

\FloatBarrier
\subsection{Estimator calibration and historical forecast diagnostics}
The calibration experiment evaluates the estimator against the conditional mean it is designed to forecast. In the independent-innovation design, the root mean bootstrap variance is 87.35 monthly basis points, compared with an estimation-error root mean square of 84.71. The ratio of mean squared error to mean estimated variance is 0.941. The estimated relative-uncertainty ranking has a mean cross-sectional correlation of 0.768 with the analytical relative-variance reference. Under feedback 0.4, the error-to-variance ratio is 0.905 and the corresponding ranking correlation is 0.754. The latter comparison uses the analytical expression as a diagnostic reference, as specified in Appendix~\ref{app:simulation}. The implied log-error dispersion is 0.414 without feedback and 0.419 with feedback. The zero-feedback value supplies the marker in Figure~\ref{fig:estimation-tradeoff}; Table~\ref{tab:calibration-error} reports its distribution across histories. True relative uncertainty in this calibration has mean coefficient of variation 75\%, between the two nonuniform settings in the controlled experiment.

Average variance accuracy coexists with conditional miscalibration. The ratio 0.941 means that mean bootstrap variance is about 6\% above mean squared error, yet nominal 95\% normal intervals cover only 92.36\% of observations. In the lowest estimated-variance decile, mean squared error is 1.89 times estimated variance and coverage is 84.53\%; in the highest decile, the ratio is 0.71 and coverage is 96.82\%. Underestimated uncertainty in one group is offset in the average by overestimated uncertainty in another, while the former group's intervals remain too narrow. Figure~\ref{fig:calibration} shows this pattern. The bootstrap variance is close to unbiased on average and ranks relative uncertainty reasonably well, so the estimates are informative but noisy. The conditional errors matter for allocation because the lowest estimated-uncertainty industries receive the weakest penalties. Appendix~\ref{app:simulation} reports the remaining coverage statistics.
\reportfigure{estimator_calibration.pdf}{Calibration of bootstrap uncertainty for conditional-mean forecasts. The first two panels compare the square root of mean estimated variance with the root mean squared conditional-mean estimation error, grouping observations by the analytical regression-variance expression or by estimated variance. The analytical expression is exact conditional on predictors in the zero-feedback design and serves as a reference under feedback 0.4. The final panel reports coverage of nominal 95\% normal and percentile intervals by estimated-variance decile. Both upper panels use monthly return basis points. Error bars use Monte Carlo standard errors clustered by simulated history.}{fig:calibration}{1}

Table~\ref{tab:forecast} presents the historical diagnostics. Forecast error exceeds that of the rolling historical-mean predictor in both periods, giving pooled predictive $R^2$ values of $-0.85\%$ and $-1.23\%$. Mean cross-sectional rank correlations are 0.0125 and 0.0035, with both intervals including zero. Meanwhile, the estimated uncertainty profile is dispersed: its mean monthly coefficient of variation is 0.779 and 0.843. Its rankings are persistent over one month, with correlations around 0.80, but twelve-month rank persistence is much lower. Forecasts therefore have higher squared error than the historical-mean benchmark despite substantial variation in estimated relative uncertainty.
\input{generated/table_forecast}

\FloatBarrier
\subsection{Net performance under two position-sizing rules}
The main economic comparison uses the common 10\% predicted-volatility target and gross cap of three (Table~\ref{tab:performance}, Panel B). Both strategies have positive sample certainty equivalents: 232 and 250 annual bp for scalar and heterogeneous allocations in 1970--2018, and 586 and 543 bp in 2019--2025. The heterogeneous-minus-scalar differences are +18 bp, with a 95\% paired interval of $[-85,88]$, and $-44$ bp, with interval $[-149,57]$. The intervals span zero in both periods. The cap binds in 51 scalar and 54 heterogeneous months out of 588 in the earlier period, and in neither strategy during the later period.

Panel A shows the direct allocations at fixed $\gamma=3$. Their certainty equivalents are negative: $-252$ and $-224$ annual bp in the earlier period, and $-109$ and $-208$ bp in the later period. Cash has zero excess-return certainty equivalent. The direct allocations earn positive average net excess returns, but their variance penalties exceed those means. The heterogeneous strategy also has higher gross exposure and turnover. Common position sizing improves absolute sample utility and reduces the magnitude of the incremental differences, while preserving their earlier positive and later negative signs.
\input{generated/table_performance}

Figure~\ref{fig:composition} shows that the estimated profile changes composition as well as scale. Mean gross-normalized composition distances are 0.536 and 0.524 in the two periods, while mean risk-direction correlations are 0.949 and 0.989. Industry weights change, but the portfolios retain closely aligned risk directions, especially in the later period. These composition measures are unchanged by positive rescaling, so they also describe the directions retained under common position sizing. Dispersion in estimated uncertainty need not produce a substantially different aggregate risk direction. Appendix~\ref{app:holdings} reports the corresponding exposure levels and penalty-selection histories.
\reportfigure{portfolio_composition.pdf}{Portfolio composition under scalar and heterogeneous uncertainty adjustments. The left panel shows the absolute distance between gross-normalized holdings; the right panel shows their correlation under the common risk matrix, defined in equation~\eqref{eq:direction-correlation}. Positive rescaling alone gives distance zero and correlation one. Both series use monthly direct-allocation portfolios at fixed risk aversion. Dates denote return realization months.}{fig:composition}{1}

Figure~\ref{fig:decomposition} decomposes the fixed-risk-aversion net utility difference in Table~\ref{tab:performance}, Panel A. In 1970--2018, the heterogeneous portfolio's gross-mean advantage contributes 243 annual bp, while higher net-return variance subtracts 175 bp and additional trading and short costs subtract 40 bp. In 2019--2025, the variance deduction exceeds the gross-mean advantage before the additional costs are included. The variance term is computed from net returns, so the decomposition includes the covariance between costs and gross returns.
\reportfigure{ce_decomposition.pdf}{Decomposition of the fixed-risk-aversion certainty equivalent difference (Table~\ref{tab:performance}, Panel A). The first four bars sum to the final net difference in each period. The gross-mean contribution is the annualized difference in average gross excess returns. Trading and short-fee contributions are the negatives of the corresponding cost differences. The final component is minus risk aversion divided by two times the annualized difference in net-return variance. All values are annual basis points.}{fig:decomposition}{1}

\FloatBarrier
\subsection{Does uncertainty-to-asset assignment add value?}
\label{sec:label-diagnostic}
Table~\ref{tab:label-placebo} compares the original assignment of uncertainty estimates to industries with 100 fixed random assignments. Each assignment preserves the monthly distribution of $\widehat h$, uses the common forecasts, risk estimates and selected penalties, and incurs its own trading costs. The statistic $D$ measures the original strategy's net certainty equivalent relative to the average utility of these random-assignment strategies.

In 1970--2018, the original assignment exceeds 99 of 100 random assignments under fixed risk aversion. Its advantage over their mean is 388.14 annual bp, with a 95\% paired interval of $[113.64,621.69]$. This provides evidence that the original assignment improves utility relative to this random benchmark under direct allocation. With common position sizing, the original assignment exceeds 88 of 100 alternatives; the mean advantage falls to 59.09 bp, with interval $[-45.43,142.32]$. The large reduction shows that the economic comparison is sensitive to position sizing. At the common risk target, the assignment advantage is estimated less precisely relative to its size.

In 2019--2025, the original assignment exceeds only 21 and 29 random assignments under the two position-sizing rules. Its mean advantages are $-80.49$ and $-42.10$ bp, with both intervals spanning zero. Thus the strong early-period result under fixed risk aversion does not persist in the later period. Table~\ref{tab:performance} separately compares the original heterogeneous rule with scalar: useful assignment information in the earlier sample coexists with a much smaller incremental gain over uniform relative uncertainty.
\input{generated/table_label_placebo}

\FloatBarrier
\subsection{Trading costs and sample eligibility}
Under the common volatility target and gross cap, raising trading cost from zero to 25 bp per dollar traded reduces the early-period net utility difference from +26 to +5 annual bp and makes the later-period difference more negative, from $-33$ to $-60$ bp (Table~\ref{tab:robustness}). Both strategies retain positive absolute certainty equivalents in every cost scenario. All six paired 12-month-block intervals for the differences include zero (Appendix~\ref{app:performance}). The early targeted gain narrows but remains positive over the tested cost range.

The direct fixed-$\gamma$ comparison is more sensitive to trading expenditure: its early-period difference falls from +62 to $-24$ bp, changing sign at the 25 bp cost rate. The later-period estimates remain negative. All cost scenarios retain the selected penalty sequence and recompute each strategy's drift, turnover and costs under the applicable position-sizing rule.

Changing the return-block length to 6 or 24 months preserves intervals spanning zero in the early period under both rules. In the later targeted comparison, the 24-month interval is $[-158.95,-2.81]$ bp; the 6- and 12-month intervals include zero. The statistical assessment of the later loss therefore depends on the block length, although no tested block length supports a positive gain. A supporting eligibility comparison requires 60 consecutive forecast-ready months under fixed $\gamma$: it changes the early estimate to 35.44 bp and leaves the later estimate unchanged, with both intervals including zero. Table~\ref{tab:robustness} distinguishes this eligibility exercise from the cost comparisons under both rules.
\input{generated/table_robustness}
\FloatBarrier
